\documentclass[11pt,a4paper]{article}
% preamble.tex — estilo común de todos los apuntes Froth.
% Cada apunte hace % preamble.tex — estilo común de todos los apuntes Froth.
% Cada apunte hace % preamble.tex — estilo común de todos los apuntes Froth.
% Cada apunte hace \input{preamble} justo después de \documentclass.
\usepackage[margin=2.4cm]{geometry}
\usepackage{xcolor}
\definecolor{litblue}{RGB}{31,79,121}
\definecolor{litgreen}{RGB}{27,94,32}
\definecolor{litorange}{RGB}{230,81,0}
\usepackage{parskip}
\usepackage{titlesec}
\titleformat{\section}{\Large\bfseries\color{litblue}}{\thesection}{0.6em}{}
\titleformat{\subsection}{\large\bfseries\color{litblue}}{\thesubsection}{0.6em}{}
\usepackage{tcolorbox}
\usepackage{fancyhdr}
\pagestyle{fancy}
\fancyhf{}
\fancyhead[L]{\small\color{litblue}Froth — Apuntes de ML}
\fancyhead[R]{\small\thepage}
\renewcommand{\headrulewidth}{0.4pt}
\usepackage[colorlinks=true,linkcolor=litblue,urlcolor=litblue]{hyperref}

% Cabecera de cada apunte: \cabecera{numero}{titulo}
\newcommand{\cabecera}[2]{%
  \begin{tcolorbox}[colback=litblue,coltext=white,colframe=litblue,arc=2mm]
    {\Large\bfseries Apunte #1 — #2}\\[3pt]
    {\small Proyecto Froth · aprender Machine Learning construyendo}
  \end{tcolorbox}\vspace{0.8em}}

% Cajas temáticas
\newtcolorbox{concepto}[1]{colback=blue!4,colframe=litblue,title={Concepto: #1},arc=1.5mm}
\newtcolorbox{practica}{colback=green!5,colframe=litgreen,title={Pruébalo tú (teclado en mano)},arc=1.5mm}
\newtcolorbox{ojo}{colback=orange!8,colframe=litorange,title={Ojo / error típico},arc=1.5mm}
\newtcolorbox{resumen}{colback=gray!8,colframe=gray!60,title={En una frase},arc=1.5mm}

\newcommand{\codigo}[1]{\texttt{#1}}
 justo después de \documentclass.
\usepackage[margin=2.4cm]{geometry}
\usepackage{xcolor}
\definecolor{litblue}{RGB}{31,79,121}
\definecolor{litgreen}{RGB}{27,94,32}
\definecolor{litorange}{RGB}{230,81,0}
\usepackage{parskip}
\usepackage{titlesec}
\titleformat{\section}{\Large\bfseries\color{litblue}}{\thesection}{0.6em}{}
\titleformat{\subsection}{\large\bfseries\color{litblue}}{\thesubsection}{0.6em}{}
\usepackage{tcolorbox}
\usepackage{fancyhdr}
\pagestyle{fancy}
\fancyhf{}
\fancyhead[L]{\small\color{litblue}Froth — Apuntes de ML}
\fancyhead[R]{\small\thepage}
\renewcommand{\headrulewidth}{0.4pt}
\usepackage[colorlinks=true,linkcolor=litblue,urlcolor=litblue]{hyperref}

% Cabecera de cada apunte: \cabecera{numero}{titulo}
\newcommand{\cabecera}[2]{%
  \begin{tcolorbox}[colback=litblue,coltext=white,colframe=litblue,arc=2mm]
    {\Large\bfseries Apunte #1 — #2}\\[3pt]
    {\small Proyecto Froth · aprender Machine Learning construyendo}
  \end{tcolorbox}\vspace{0.8em}}

% Cajas temáticas
\newtcolorbox{concepto}[1]{colback=blue!4,colframe=litblue,title={Concepto: #1},arc=1.5mm}
\newtcolorbox{practica}{colback=green!5,colframe=litgreen,title={Pruébalo tú (teclado en mano)},arc=1.5mm}
\newtcolorbox{ojo}{colback=orange!8,colframe=litorange,title={Ojo / error típico},arc=1.5mm}
\newtcolorbox{resumen}{colback=gray!8,colframe=gray!60,title={En una frase},arc=1.5mm}

\newcommand{\codigo}[1]{\texttt{#1}}
 justo después de \documentclass.
\usepackage[margin=2.4cm]{geometry}
\usepackage{xcolor}
\definecolor{litblue}{RGB}{31,79,121}
\definecolor{litgreen}{RGB}{27,94,32}
\definecolor{litorange}{RGB}{230,81,0}
\usepackage{parskip}
\usepackage{titlesec}
\titleformat{\section}{\Large\bfseries\color{litblue}}{\thesection}{0.6em}{}
\titleformat{\subsection}{\large\bfseries\color{litblue}}{\thesubsection}{0.6em}{}
\usepackage{tcolorbox}
\usepackage{fancyhdr}
\pagestyle{fancy}
\fancyhf{}
\fancyhead[L]{\small\color{litblue}Froth — Apuntes de ML}
\fancyhead[R]{\small\thepage}
\renewcommand{\headrulewidth}{0.4pt}
\usepackage[colorlinks=true,linkcolor=litblue,urlcolor=litblue]{hyperref}

% Cabecera de cada apunte: \cabecera{numero}{titulo}
\newcommand{\cabecera}[2]{%
  \begin{tcolorbox}[colback=litblue,coltext=white,colframe=litblue,arc=2mm]
    {\Large\bfseries Apunte #1 — #2}\\[3pt]
    {\small Proyecto Froth · aprender Machine Learning construyendo}
  \end{tcolorbox}\vspace{0.8em}}

% Cajas temáticas
\newtcolorbox{concepto}[1]{colback=blue!4,colframe=litblue,title={Concepto: #1},arc=1.5mm}
\newtcolorbox{practica}{colback=green!5,colframe=litgreen,title={Pruébalo tú (teclado en mano)},arc=1.5mm}
\newtcolorbox{ojo}{colback=orange!8,colframe=litorange,title={Ojo / error típico},arc=1.5mm}
\newtcolorbox{resumen}{colback=gray!8,colframe=gray!60,title={En una frase},arc=1.5mm}

\newcommand{\codigo}[1]{\texttt{#1}}

\begin{document}
\cabecera{25}{De la teoría al texto: find\_gaps() y el borrador de review}

\section{Qué construimos (5.1 + 5.2)}

Las dos mitades del diferenciador, ya funcionando en \codigo{froth/gaps.py}:
\begin{enumerate}
  \item \codigo{find\_gaps()}: convierte los 4 tipos del Apunte 24 en señales medidas
        $\rightarrow$ tabla rankeada con evidencia (\codigo{2\_Datos/processed/gaps.parquet}).
  \item \codigo{draft\_review()}: redacta el andamio del literature review
        (\codigo{3\_Resultados/review\_draft.md}) con \textbf{citas 100\% trazables}.
\end{enumerate}
Se corren juntas: \codigo{python -m froth.gaps}.

\section{Los gaps de TU corpus (resultados reales)}

\begin{center}
\begin{tabular}{lp{6.2cm}l}
\textbf{Tipo} & \textbf{Candidato top} & \textbf{Evidencia} \\
\hline
Silo & commodities/críticos $\leftrightarrow$ flotación & sim 0.69, \textbf{1} cita cruzada \\
Silo & flotación $\leftrightarrow$ geología (\emph{tu tesis}) & sim 0.70, 3 citas cruzadas \\
Puente escaso & commodities $\leftrightarrow$ geología & solo 8 puentes pese a sim 0.67 \\
Dormido & (ninguno: campo caliente, papers 2025-26) & la herramienta lo dice honesto \\
\end{tabular}
\end{center}

Lectura fina del silo campeón: la literatura de \emph{criticidad/cadenas de suministro} y la
de \emph{procesamiento} casi no se citan. Conectarlas (``mi flotación importa para la
seguridad de suministro de Li/Ta'') es munición directa para introducciones y propuestas.

\begin{ojo}
La señal más débil es la \textbf{zona rala}: mide \emph{ausencias}, y un vacío no explica su
porqué (¿mina o desierto?) — además de heredar las distorsiones del mapa 2D (Apunte 12). Por
eso va etiquetada ``experimental'' en el código y en el draft. Regla general: desconfía más
de las métricas sobre lo que falta que de las métricas sobre lo que existe.
\end{ojo}

\section{El borrador: por qué una plantilla y no un LLM}

\begin{concepto}{citas trazables o nada}
El draft se genera con una \textbf{plantilla determinista}: cada sección nace de la
estructura MEDIDA (secciones $=$ clusters; imprescindibles $=$ más citados + hubs locales;
novedades $=$ más recientes; huecos $=$ gaps con sus números) y cada afirmación apunta a
papers reales con referencia numerada [n] $\rightarrow$ lista final con autor/año/título/DOI.
Un LLM redactaría más bonito... y tarde o temprano \emph{inventaría} una referencia — el
pecado capital académico. Orden correcto: primero el andamio confiable, después (opcional)
un LLM que pula la prosa \emph{verificando} cada cita contra el corpus.
\end{concepto}

\section{Frente a la competencia (otro vecindario)}

Aquí los rivales no son los mapeadores sino Elicit/SciSpace/Consensus (LLM que sintetiza
\emph{sin} conocer la estructura del corpus), Scholarcy (resúmenes 1-a-1) y ChatGPT a pelo
(referencias alucinadas). \textbf{Nadie une mapa y texto}: nuestro draft dice ``estos dos
subtemas comparten 0.69 de similitud y una sola cita cruzada'' porque lo midió. Límite
honesto: trabajamos con abstracts (Elicit extrae del full text) $\rightarrow$ lo nuestro es
un \emph{andamio estructurado}, no prosa final.

\begin{resumen}
find\_gaps() produce candidatos rankeados con evidencia (silos al frente: commodities-
flotación 0.69/1 cita; tu tesis reconfirmada 0.70/3); draft\_review() los convierte en un
andamio de review con 26 referencias trazables, sin LLM. Nadie más une mapa y texto.
\end{resumen}

\end{document}
